\chapter{Giới thiệu}
\label{chap:introduction}

\section{Bối cảnh và Động lực}
\label{sec:intro-motivation}

Trong chương trình đào tạo ngành Cơ Kỹ Thuật tại Trường Đại học Bách Khoa TP.HCM,
sinh viên được yêu cầu thực hiện các bài toán thiết kế hệ thống dẫn động cơ khí như:
tính chọn động cơ điện, thiết kế bộ truyền xích và bộ truyền bánh răng.
Các bài toán này đòi hỏi thực hiện tuần tự nhiều bước tính toán phức tạp,
tra cứu nhiều bảng số liệu tiêu chuẩn và áp dụng hàng chục công thức từ giáo trình~\cite{trinh2006tinh,Loc2023GiaoTrinh,Loc2020ThietKe}.

Thực tế hiện nay, toàn bộ quá trình tính toán được thực hiện thủ công bằng bút giấy
hoặc bảng tính Excel rời rạc, dẫn đến một số vấn đề điển hình:
\begin{itemize}
  \item Dễ sai số do tra bảng thủ công hoặc nhầm công thức giữa các bước.
  \item Khi thay đổi một thông số đầu vào (ví dụ: thay động cơ), phải tính lại toàn bộ chuỗi.
  \item Không có công cụ kiểm tra chéo tức thời giữa các module (động cơ -- xích -- bánh răng).
\end{itemize}

Xuất phát từ thực tiễn đó, đồ án đa ngành này được thực hiện nhằm xây dựng một
ứng dụng di động giúp số hóa toàn bộ quy trình tính toán thiết kế,
kết hợp kiến thức cơ kỹ thuật với phương pháp luận kỹ thuật phần mềm hiện đại.

\section{Mục tiêu Đề tài}
\label{sec:intro-objectives}

Đề tài hướng đến các mục tiêu cụ thể sau:

\begin{enumerate}
    \item Xây dựng ứng dụng di động hỗ trợ tính toán 3 module theo đề 2 phương án 2 đồ án thiết kế cơ kỹ thuật (hệ dẫn động thùng trộn):
        \begin{itemize}
            \item \textbf{Module 1}: Tính chọn động cơ điện và phân phối tỉ số truyền.
            \item \textbf{Module 2}: Tính thiết kế và kiểm nghiệm bộ truyền xích ống con lăn.
            \item \textbf{Module 3}: Tính thiết kế và kiểm nghiệm bộ truyền bánh răng
                (côn răng thẳng -- cấp nhanh; trụ răng thẳng -- cấp chậm).
        \end{itemize}
    \item Áp dụng ít nhất 1 Design Pattern trong thiết kế phần mềm theo yêu cầu hướng Công nghệ Phần mềm của đồ án đa ngành.
    \item Đảm bảo kết quả tính toán của ứng dụng nhất quán với kết quả tính tay theo phương án 2 đang được triển khai của sinh viên Cơ kỹ thuật \cite{LeNhatTruong2026}.
\end{enumerate}

\section{Phạm vi Thực hiện}
\label{sec:intro-scope}

\textbf{Trong phạm vi:}
\begin{itemize}
    \item Tính toán cho hệ thống dẫn trục được đề xuất trong Đề số 2 của Đồ án Thiết kế ngành Cơ kỹ thuật học kỳ 252
    \item Tính toán ba module kể trên, lấy bộ số liệu Phương án~2 làm kiểm thử:
            công suất $P = 7{,}5$~kW, số vòng quay $n = 75$~v/p, thời gian phục vụ $L = 10$~năm.
    \item Các biến tùy chỉnh trong các giáo trình \cite{trinh2006tinh, Loc2023GiaoTrinh, Loc2020ThietKe} được cố định theo đồ án thực tế\cite{LeNhatTruong2026}
    \item Ứng dụng di động đa nền tảng (Android và iOS) xây dựng bằng React Native cùng với Backend được xây dựng bởi Java Spring + PostGreSql.
    \item Cơ sở dữ liệu cục bộ lưu bảng tra cứu động cơ tiêu chuẩn và bảng vật liệu.
    \item Lưu/tải lại phiên tính toán trên thiết bị.
\end{itemize}

\textbf{Ngoài phạm vi:}
\begin{itemize}
    \item Hỗ trợ tính toán cho mọi loại thiết kế hệ thống dẫn trục
    \item Hỗ trợ ngoài ba module được liệt kê ở trên.
    \item Hỗ trợ người dùng điều chỉnh các biến tùy chỉnh
\end{itemize}

\section{Cấu trúc Báo cáo}
\label{sec:intro-structure}

Báo cáo được tổ chức thành 7 chương:

\begin{itemize}
  \item \textbf{Chương~\ref{chap:introduction} -- Giới thiệu}: Bối cảnh, mục tiêu, phạm vi đề tài.
  \item \textbf{Chương~\ref{chap:overview} -- Tổng quan}: Kiến thức cơ sở cơ kỹ thuật,
        công nghệ sử dụng và các công trình liên quan.
  \item \textbf{Chương~\ref{chap:requirements} -- Phân tích Yêu cầu}: Yêu cầu chức năng,
        phi chức năng, use case diagram và đặc tả use case.
  \item \textbf{Chương~\ref{chap:design} -- Thiết kế Hệ thống}: Kiến trúc phần mềm,
        áp dụng Design Pattern, thiết kế CSDL và giao diện người dùng.
  \item \textbf{Chương~\ref{chap:implementation} -- Hiện thực}: Chi tiết cài đặt
        từng module và các pattern đã thiết kế.
  \item \textbf{Chương~\ref{chap:testing} -- Kiểm thử}: Chiến lược, test case
        và kết quả kiểm thử.
  \item \textbf{Chương~\ref{chap:conclusion} -- Kết luận}: Tổng kết kết quả,
        hạn chế và hướng phát triển.
\end{itemize}
